% Generated draft: requires mathematical and source review.
\documentclass[11pt]{article}
\usepackage[margin=25mm]{geometry}
\usepackage{fontspec}
\setmainfont{Latin Modern Roman}
\usepackage{amsmath,amssymb,mathtools,unicode-math}
\setmathfont{Latin Modern Math}
\usepackage{enumitem,needspace,xurl,hyperref}
\setlength{\parindent}{0pt}
\setlength{\parskip}{5pt}
\setlength{\emergencystretch}{3em}
\newcommand{\N}{\mathbb{N}}
\newcommand{\Z}{\mathbb{Z}}
\newcommand{\Q}{\mathbb{Q}}
\newcommand{\R}{\mathbb{R}}
\newcommand{\C}{\mathbb{C}}
\newcommand{\F}{\mathbb{F}}
\newcommand{\A}{\mathbb{A}}
\newcommand{\PP}{\mathbb{P}}
\newcommand{\E}{\mathbb{E}}
\newcommand{\eps}{\varepsilon}
\newcommand{\dd}{\,\mathrm{d}}
\begin{document}
\section*{UnsolvedMath ID: 2306039}
\textbf{Source problem code:} catalogue entry
\subsection{Definitions and mathematical statement}
Let $S$ be the class of injective holomorphic functions $f:\mathbb D\to\C$, $\mathbb D=\{|z|<1\}$, normalized by $f(0)=0$, $f'(0)=1$. Define
\[
h(z)=\bigl(f(z^2)\bigr)^{1/2}
=z+\sum_{k=2}^\infty c_{2k-1}z^{2k-1}.
\]
The square root is the holomorphic branch with $h'(0)=1$: explicitly $h(z)=z\sqrt{f(z^2)/z^2}$ with the square root equal to $1$ at $0$. The quotient extends holomorphically and is zero-free on the simply connected disc, so this normalization fixes the branch. Set $c_1=1$.

Robertson's conjecture is that for every $f\in S$ and integer $n\geq1$,
\[
\sum_{k=1}^n|c_{2k-1}|^2
=1+|c_3|^2+\cdots+|c_{2n-1}|^2\leq n.
\]
At $n=1$ the statement is the normalized equality $1\leq1$.

The source says this was known for starlike $f$ and asks in particular whether it holds for $f$ with real Taylor coefficients, or for close-to-convex $f$. Starlike means that $f(\mathbb D)$ is star-shaped about $0$. A standard analytic close-to-convex convention is that there exist a normalized starlike function $g$ and a real $\theta$ with $\operatorname{Re}(e^{i\theta}zf'(z)/g(z))>0$ throughout the disc, using the removable value at $0$. Real coefficients means the normalized Taylor coefficients of $f$ all lie in $\R$.

The supplied research summary reports Robertson's inequality for all $f\in S$ as a consequence of de Branges's inequalities in his proof of the Bieberbach conjecture. That general result answers both requested subclasses affirmatively. No equality classification beyond the normalized $n=1$ observation is added.
\subsection{Short English statement}
The normalized odd square-root transform of a univalent function is conjectured to have sum of squared first $n$ odd coefficients at most $n$. The supplied record reports the full Robertson conjecture proved through de Branges's work, covering the real-coefficient and close-to-convex questions.
\subsection{Sources}
\begin{itemize}
\item Local UnsolvedMath registry, permanent ID 2306039, exact source code catalogue entry, ``Research Problems in Function Theory --- Problem 6.39''; status solved, classification SOLVED-IN-LITERATURE.
\item Hayman, \emph{Research Problems in Function Theory} (2018), Problem 6.39, attributed to P. L. Duren, \url{https://arxiv.org/abs/1809.07200}. The record identifies extraction from public arXiv source TeX, accessed July 29, 2026; the 2018 update was retained for status triage.
\item Robertson's conjecture, with Sheil-Small cited in the source without a complete bibliographic entry; the source records the starlike case as known.
\item De Branges's inequalities and proof of the Bieberbach conjecture, cited by the supplied research summary as proving Robertson's conjecture for all $n$. This record gives no full bibliographic entry or theorem number.
\end{itemize}
\end{document}
